% Титульный лист для загрузки в Антиплагиат
% Не нумеруется, располагается перед содержанием

\begin{titlepage}

\begin{center}
{\LARGE\bfseries Устинов Андрей Эдуардович}
\end{center}

\vspace{1cm}

\begin{center}
{\LARGE\bfseries Исследование когнитивных искажений больших языковых моделей}
\end{center}

\vspace{1.5cm}

\noindent\textbf{ЗАДАНИЕ НА ВЫПУСКНУЮ КВАЛИФИКАЦИОННУЮ РАБОТУ}

\vspace{0.5cm}

\begin{enumerate}
  \item Провести анализ существующих подходов к исследованию когнитивных предубеждений в больших языковых моделях и выявить методологические ограничения предшествующих работ.
  \item Разработать экспериментальный дизайн для трёх психологических парадигм (задача Линды, задача выбора карточек Уэйсона, задача открытия правила 2-4-6) с контролем контаминации обучающих данных.
  \item Сформировать оригинальные датасеты для каждого эксперимента, обеспечивающие структурное единообразие стимулов и деконтаминацию относительно обучающих корпусов тестируемых моделей.
  \item Разработать систему метрик для количественной оценки когнитивных предубеждений в языковых моделях.
  \item Провести сравнительное экспериментальное исследование на наборе современных языковых моделей и выполнить статистический анализ полученных результатов.
  \item Выполнить сравнительный анализ профилей когнитивных предубеждений по трём парадигмам и проверить сформулированные исследовательские гипотезы.
  \item Сформулировать рекомендации по выбору и оценке LLM для приложений, требующих нормативно корректного рассуждения.
\end{enumerate}

\vspace{1cm}

\noindent\textbf{АННОТАЦИЯ}

\vspace{0.5cm}

\noindent\textbf{Цель работы}~--- систематически измерить и сравнить когнитивные предубеждения в современных больших языковых моделях (LLM) на трёх психологических парадигмах: вероятностное суждение (задача Линды), условная логика (задача выбора карточек Уэйсона) и индуктивное открытие правил (задача 2-4-6). Объект исследования~--- LLM как системы рассуждения, предмет~--- когнитивные предубеждения в их ответах. Методы: парадигмы когнитивной психологии, адаптированные для LLM, точный тест МакНемара, процедура Бенджамини-Хохберга, метрика IOU на основе Монте-Карло. Научная новизна: впервые три парадигмы исследованы в едином дизайне на девяти моделях; введена метрика Severity of Bias; разработан датасет из 130 виньеток с факториальными подстановками; введено условие «фантастических социальных контрактов».

\vspace{0.5cm}

\noindent\textbf{Задачи:} анализ существующих подходов и выявление ограничений предшествующих работ; разработка экспериментального дизайна с контролем контаминации; формирование оригинальных датасетов; разработка системы метрик; проведение экспериментов на девяти моделях; сравнительный анализ профилей предубеждений; формулирование практических рекомендаций по выбору LLM.

\vspace{0.5cm}

\noindent\textbf{Краткие результаты:} ошибка конъюнкции выявлена у восьми из девяти моделей; предубеждение подтверждения воспроизведено у всех девяти; в индуктивной задаче CTR $\geq 0{,}629$ у всех non-reasoning моделей; масштаб модели не является монотонным предиктором устойчивости; reasoning-режим обеспечивает качественный скачок, недостижимый через prompt engineering; деонтическая структура фасилитирует вывод независимо от знакомости содержания.

\end{titlepage}
